\chapter{Sustainability analysis} 
Considerations regarding sustainability has become an important element of modern aerospace engineering. Therefore for this project it was very important to design for the sustainable construction and operation of the aircraft. Sustainability should always be a consideration for any engineer involved with this project. 
\subsection*{Construction}
Historically the construction of an aircraft have been a source of large emissions. The requirements for exotic materials with a high degree of processing leads to large energy usage and generation of waste-products. This has to be taken into consideration for the choice of material and design of the aircraft. There should be an emphasis on minimizing the use of materiel by the construction of efficient structures. Another important consideration is that construction is a one in a lifetime emission for the aircraft. This means that a weight saving measure that  results in higher production emission might be worth it for the longer terms emission reduction due to the weight saving.         
\subsection*{Operation}
Since it was specified that the aircraft should use the Vulcan turbine engine the aircraft will always generate emissions while in use. It was also specified that this turbine should be powered by SAF (sustainable aviation fuel) meaning the carbon-emissions released onto the atmosphere would originate from a capture process. This would lead to net-zero carbon emissions. Though this does not mitigate the generation of NOX typically contributed to turbine jets. It is therefore also relevant to design for modes of operation and implement systems that would decrease the amount of NOX.      
\subsection*{End-of-life}
The aircraft is designed to have a long lifetime but it is also important to design for the point when the aircraft needs to be scrapped. Recyclability should be prioritized for all of the components of the aircraft. This is to be considered in the operation but especially in the design and construction of the aircraft. The material choice and bonding methods will influence how difficult (or whether it is even possible) to recycle the used material. 

\subsection*{Life Cycle Assessment}
To quantify the material input and output and the subsequent consequences of these the Life Cycle Assessment model is used to evaluate the more detailed concepts on later stages of the design process. Since a full Life Cycle Assessment is very time consuming it is relevant to use the streamlined Life Cycle Assessment. This is a simplified model which focuses on the differences between two designs and can be used to quickly evaluate more different designs. (Insert citation)   

\subsection*{Environmental benefits of product}
Another important aspect to consider is the overall environmental impact of the research enabled by the system. An emphasis was placed on how the system could enable a more detailed understanding of the operation of high-aspects ratio wings. 
%The implementation of such wings could lead to a potential reduction in fuel-usage of around 10\%   
The utilization of such technologies could potentially lead to major reductions in aviation-emissions (by some estimates up to 10\% Insert citation). The major benefit of this product is the rapid testing and iteration of potential fuel-saving designs that it enables.    